\section{はじめに}
大喜利は，出題された題目に対して機知に富んだ回答を返す言葉遊びであり，ユーモアの生成と理解を扱う研究の題材として用いられている．
大喜利投稿サイト「ボケて」\cite{ref_bokete}では，投稿者がお題となる画像を掲示し，他の利用者がその画像に対する面白い一言を投稿する．
ボケてに蓄積された画像と回答の組は，画像を入力として面白い文を生成するモデルの学習や，ユーモアの分析に利用できる大規模な資源である．
本稿では，ボケてから収集したデータベースを，学習や分析に用いやすい形に整えるため，全てのお題画像に種類を表すラベルを付与する作業を扱う．
以下では，データにラベルを付与する作業をアノテーションと呼ぶ．

お題画像には，現実の写真のように画像を見て一言を付ければよいものだけでなく，回答の形を制限するものがある．
例えば，画像中の空欄に入る語を答える穴埋め，画像の隅に書かれたひらがな1文字から始まる文を作るカルタ，画像に書かれた問いや指示に答えるお題文がある．
また，アニメのスクリーンショットや漫画のコマのように，作品の知識を前提とする画像もある．
これらの画像からボケを生成するには，空欄の位置の検出，ひらがなの読み取り，文字の読み取りなど，種類ごとに異なる処理が必要になる．
したがって，生成モデルの学習やデータの分析では，各お題画像がどの種類に属するかを事前に知る必要がある．

アノテーションの対象は約682万件のお題画像であり，人手で全てにラベルを付けることは現実的ではない．
近年の視覚言語モデル (Vision-Language Model，VLM) は，画像と指示文を入力として，指示に従った分類を学習なしで行える．
しかし，有料の API で約682万件を処理すると費用が大きく，ローカルの GPU で VLM を全件に適用しても約70時間以上を要する．
さらに，ラベルの定義を見直すたびに全件を処理し直す必要がある．
そこで本稿では，有料の API を用いず，単一の GPU で全件を現実的な時間で処理できるアノテーションの方法を確立することを目的とする．

本研究は，方法の比較，ラベルの定義の改訂，評価の厳密化を繰り返す試行錯誤の過程を経て，最終の方法に到達した．
本稿では，最終の方法だけでなく，その過程で得られた知見を，以下の観点からまとめる．
\begin{itemize}
    \item ゼロショット分類，メタデータによる弱教師あり学習，VLM による直接の分類，VLM のラベルを高速な分類器に移す知識蒸留を比較し，蒸留が精度と速度の両面で最良であることを示す．
    \item 強い手法に共通する誤りの多くがラベルの定義の境界にあることを示し，ボケの生成に必要な処理から逆算してラベルの定義を改訂する過程と，その各段階での判断の根拠を示す．
    \item VLM への指示の方法について，質問を分割する方法，説明を書き足す方法，別の短い質問で確かめる方法などの成否を比較し，有効であった方法とそうでなかった方法の傾向を示す．
    \item 出現率の低い区分を扱うための学習データの追加，重み付きの損失，複数の分類器の確率の平均を比較し，最終の構成を決める．
    \item 最終の構成を全データに適用し，ラベルの分布と投稿年ごとの推移を示す．
\end{itemize}

本稿の構成は次のとおりである．
\ref{sec_data}章で対象のデータと本稿で用いる用語を定義し，\ref{sec_related}章で関連する技術を整理する．
\ref{sec_eval}章で評価用データと評価指標を定義する．
\ref{sec_selection}章で手法の選定の実験を述べ，\ref{sec_label}章でラベルの定義の改訂の経緯を述べる．
\ref{sec_teacher}章と\ref{sec_student}章で，教師となる VLM への指示の方法と，生徒となる分類器の構成をそれぞれ検討する．
\ref{sec_final}章で最終の構成とその精度を示し，\ref{sec_full}章で全データへの適用の結果を示す．
\ref{sec_discussion}章で得られた知見と限界を考察し，\ref{sec_conclusion}章でまとめる．
